\question[15]

%\begin{minipage}[t]{0.4\textwidth}
\begin{multicols}{2}
\flushleft
The table at right (Sharp et al. 2005) shows how \textit{E. coli} uses 32 codons from 40 highly expressed genes. The tabled numbers are the ratio of actual use to expected use if all codons were used equally. A ratio of 1 means the codon is used exactly as often as expected (\textit{e.g.}, Met, in bold). Explain whether the data in the table provide evidence for natural selection at the genetic level and, if so, how. If not, explain why not. Use 1--2 specific codons from the table to support your answer.
%\end{minipage}\hspace{1em}\begin{minipage}[t]{0.5\textwidth}
\columnbreak
\begin{tabular}{@{}llrcllr@{}}
	\toprule
	\multicolumn{7}{l}{Codon use in \textit{E. coli}.}\\
	\midrule
	Phe & UUU & 0.45 &  & Ser & UCU & 2.54\\
Phe & UUC & 1.55 &  & Ser & UCC & 1.52\\
Leu & UUA & 0.14 &  & Ser & UCA & 0.19\\
Leu & UUG & 0.25 &  & Ser & UCG & 0.06\\
Leu & CUU & 0.30 &  & Pro & CCU & 0.59\\
Leu & CUC & 0.23 &  & Pro & CCC & 0.05\\
Leu & CUA & 0.01 &  & Pro & CCA & 0.52\\
Leu & CUG & 5.07 &  & Pro & CCG & 2.85\\
Ile & AUU & 0.72 &  & Thr & ACU & 1.89\\
Ile & AUC & 2.27 &  & Thr & ACC & 1.75\\
Ile & AUA & 0.01 &  & Thr & ACA & 0.19\\
\textbf{Met} & \textbf{AUG} & \textbf{1.00} &  & Thr & ACG & 0.17\\
Val & GUU & 2.07 &  & Ala & GCU & 1.83\\
Val & GUC & 0.30 &  & Ala & GCC & 0.32\\
Val & GUA & 1.14 &  & Ala & GCA & 1.09\\
Val & GUG & 0.48 &  & Ala & GCG & 0.76\\
	\bottomrule
	\end{tabular}
\end{multicols}
%	\end{minipage}
	\begin{solution}
	Codon bias is present.
	\end{solution}

